%===============================================================
% 07_developer.tex - 開発者向けドキュメント・設計詳細
%===============================================================
\section{開発者向け設計詳細と拡張性}

本システムは、将来の機能追加やLLMプロバイダの変更、UIフレームワークの刷新に耐えうるよう、疎結合なモジュール設計を徹底しています。

\subsection{疎結合設計の原則}

本システムにおける各モジュールの結合は、HTTP API（REST / SSE）または明確なPythonクラスメソッドの引数/戻り値に限定されています。これにより、以下の変更が生じた場合でも、他のコンポーネントに影響を与えることなく改修が可能です。

\begin{enumerate}
  \item \textbf{LLM（AI）プロバイダの変更}: 
    Ollamaから外部API（OpenAI, Anthropicなど）に移行する場合でも、\file{server/modules/llm\_client.py} 内の \texttt{OllamaClient} クラス（またはこれに代わる共通インターフェース）の内部処理を書き換えるだけで、FastAPI のルーティングやフロントエンド側には一切の変更が発生しません。
  \item \textbf{フロントエンドの刷新}:
    現在の Vite + React から他のフレームワーク（Next.js や Vue.js など）へ移行、あるいはネイティブアプリ（iOS / Android）を作成する場合でも、バックエンドの OpenAPI (FastAPI) が提供するエンドポイント群（\port{8000}）は同一であるため、バックエンドコードの変更は一切不要です。
  \item \textbf{MIDI変換エンジンの変更}:
    現在は C++ で書かれた実行バイナリ \cmd{abc2midi} を Node.js でラップして処理していますが、将来的に Python のネイティブライブラリによる処理へ移行する場合、\file{server/routes/midi\_convert.py} を書き換えて Node.js サーバーへのプロキシを解除するだけで、フロントエンド側の変換インターフェースは維持されます。
\end{enumerate}

\subsection{モジュール依存関係図}

フロントエンドおよびバックエンドにおけるコードモジュール同士の依存関係を図\ref{fig:dependency}に示します。

\begin{figure}[H]
  \centering
  \begin{tikzpicture}[
    scale=0.75, every node/.style={transform shape},
    node distance=1.2cm and 1.5cm,
    % スタイル定義
    fe/.style={draw=node-ui, fill=node-ui!10, rounded corners=4pt, minimum width=2.8cm, minimum height=0.7cm, font=\small\ttfamily, align=center},
    be/.style={draw=node-api, fill=node-api!10, rounded corners=4pt, minimum width=2.8cm, minimum height=0.7cm, font=\small\ttfamily, align=center},
    ext/.style={draw=node-midi, fill=node-midi!10, rounded corners=4pt, minimum width=2.8cm, minimum height=0.7cm, font=\small\ttfamily, align=center},
    ol/.style={draw=node-ollama, fill=node-ollama!10, rounded corners=4pt, minimum width=2.8cm, minimum height=0.7cm, font=\small\ttfamily, align=center},
    dep/.style={-{Latex[length=6pt]}, thick, draw=brand-purple!70}
  ]

  % Frontend Nodes
  \node[fe] (app) {App.jsx (Main)};
  \node[fe, below left=1cm and 0.5cm of app] (comp) {ComposeTab};
  \node[fe, below=1cm of app] (harm) {HarmonizeTab};
  \node[fe, below right=1cm and 0.5cm of app] (midiconv) {MidiConvertTab};
  \node[fe, right=2cm of app] (chat) {ChatTab};

  \draw[dep] (app) -- (comp);
  \draw[dep] (app) -- (harm);
  \draw[dep] (app) -- (midiconv);
  \draw[dep] (app) -- (chat);

  % Backend Nodes (FastAPI)
  \node[be, below=3.5cm of app] (main) {main.py (FastAPI)};
  \node[be, below left=1cm and 0.5cm of main] (rcompose) {routes/compose.py};
  \node[be, below=1.2cm of main] (rharmonize) {routes/harmonize.py};
  \node[be, below right=1cm and 0.5cm of main] (rmidi) {routes/midi\_convert.py};
  \node[be, right=2.2cm of main] (rchat) {routes/chat.py};

  \draw[dep] (main) -- (rcompose);
  \draw[dep] (main) -- (rharmonize);
  \draw[dep] (main) -- (rmidi);
  \draw[dep] (main) -- (rchat);

  % Shared Backend modules
  \node[be, below=1.2cm of rharmonize] (llmclient) {modules/llm\_client.py};
  \node[be, left=1.5cm of llmclient] (synthesizer) {abc\_synthesizer.py};
  \node[be, below=0.8cm of llmclient] (config) {modules/config.py};
  
  \draw[dep] (rcompose) -- (llmclient);
  \draw[dep] (rharmonize) -- (llmclient);
  \draw[dep] (rchat) -- (llmclient);
  
  \draw[dep] (rcompose) -- (synthesizer);
  \draw[dep] (rharmonize) -- (synthesizer);
  \draw[dep] (llmclient) -- (config);

  % External Connectors
  \node[ext, right=1.5cm of rmidi] (node-srv) {midi-server/server.js};
  \node[ol, below=1.2cm of rchat] (ollama-api) {Ollama API};

  \draw[dep] (rmidi) -- (node-srv);
  \draw[dep] (llmclient) -- (ollama-api);

  % Frontend to Backend (HTTP Connectors)
  \draw[dep, dotted, brand-cyan] (comp.south) .. controls +(0,-1) and +(0,1) .. (rcompose.north);
  \draw[dep, dotted, brand-cyan] (harm.south) -- (rharmonize.north);
  \draw[dep, dotted, brand-cyan] (midiconv.south) .. controls +(0,-1) and +(0,1) .. (rmidi.north);
  \draw[dep, dotted, brand-cyan] (chat.south) .. controls +(0,-1) and +(0,1) .. (rchat.north);

\end{tikzpicture}
\caption{フロントエンド・バックエンドにおけるモジュール依存関係}
\label{fig:dependency}
\end{figure}
